\documentclass[french,a4paper]{article}

\usepackage[utf8]{inputenc}
\usepackage[T1]{fontenc}
\usepackage{lmodern}
\usepackage{geometry}
\usepackage{graphicx}
\usepackage{babel}
\usepackage{amsmath}
\usepackage{amsthm}
\usepackage{amssymb}
\usepackage{hyperref}
\usepackage{stmaryrd}
\usepackage{enumitem}
\usepackage{tcolorbox}
\usepackage{dsfont}
\usepackage{multirow}
\usepackage{etoc}
\usepackage{titlesec}
\usepackage{esint}
\usepackage{cancel}
\usepackage{mathdots}

\setlist[itemize]{label=$\bullet$}


\renewcommand{\P}{\mathbb{P}}
\newcommand{\N}{\mathbb{N}}
\newcommand{\Z}{\mathbb{Z}}
\newcommand{\Q}{\mathbb{Q}}
\newcommand{\R}{\mathbb{R}}
\newcommand{\C}{\mathbb{C}}
\newcommand{\K}{\mathbb{K}}
\renewcommand{\L}{\mathbb{L}}
\newcommand{\U}{\mathbb{U}}
\newcommand{\st}{^*}
\newcommand{\tq}{\middle|}
\newcommand{\inv}{^{-1}}
\newcommand{\E}{\mathbb{E}}
\newcommand{\F}{\mathbb{F}}
\newcommand{\V}{\mathbb{V}}
\renewcommand{\ker}{\text{Ker}}
\newcommand{\im}{\text{Im}}
\newcommand{\card}{\text{Card}}
\newcommand{\Tr}{\text{Tr}}

\theoremstyle{plain}
\newtheorem{thm}{Théorème}
\newtheorem*{ind}{Indication}

\theoremstyle{definition}
\newtheorem{dfn}{Definition}

\theoremstyle{remark}
\newtheorem*{rmq}{Remarque}
\newtheorem*{app}{Application}

\newtcolorbox[auto counter]{ex}[2][]{colback=white, colframe=green!50!white, coltitle=black, fonttitle=\bfseries, title={Exercice~\thetcbcounter~: #2}, after title={\hfill #1}}

\title{TD Réduction}

\begin{document}

\maketitle

Dans toute cette feuille, $\K$ désigne un sur-corps de $\Q$.

\section{Éléments propres}

\begin{ex}[Moulin.01]{}
	Déterminer les valeurs propres de l'endomorphisme $u$ de $\R[X]$ défini par : $$\forall P\in\R[X],\ u(P)=(X^2-1)P''+(2X+1)P'$$
\end{ex}

\begin{ex}[Moulin.02]{}
	Déterminer les valeurs propres et les vecteurs de la matrice : $$\begin{pmatrix}
		1&1&\cdots&1\\
		1&0&\cdots&0\\
		\vdots&\vdots&&\vdots\\
		1&0&\cdots&0
	\end{pmatrix}$$
\end{ex}

\begin{ex}[Moulin.03]{$\star$}
	Déterminer les valeurs propres et les vecteurs propres de l'endomorphisme $u$ de $\mathcal{C}^0(\R,\C)$ défini par : $$u(f)(x)=\int_{0}^{2\pi}\sin(x-t)f(t)\ dt$$
\end{ex}

\begin{ex}[Moulin.04]{$\star$}
	Soit $u$ un endomorphisme d'un espace vectoriel de dimension finie.
	\begin{enumerate}
		\item Montrer que $u^T:\varphi\mapsto\varphi\circ u$ est un endomorphisme de $E\st$.
		\item Montrer que $u^T$ et $u$ ont les même valeurs propres.
		\item Montrer que $u$ laisse stable un hyperplan de $E$ si, et seulement si, son polynôme caractéristique admet une racine.
		\item Déterminer la matrice de $u^T$ dans la base duale d'une base $\mathcal{B}$ donnée de $E$.
	\end{enumerate}
\end{ex}

\begin{ex}[Moulin.05]{$\star$}
	Pour $A\in\C_{n-1}[X]$, on définit $\Phi:\C_{n-1}[X]\to\C_{n-1}[X]$ en posant $\Phi(P)$ le reste de la division euclidienne de $AP$ par $X^n-1$. Déterminer les éléments propres de $\Phi$. \\
	Application : diagonalisation de la matrice circulante $$\begin{pmatrix}
		a_0&a_{n-1}&\cdots&a_1\\
		a_1&a_0&\ddots&\vdots\\
		\vdots&\vdots&\ddots&a_{n-1}\\
		a_{n-1}&a_{n-2}&\cdots&a_0
	\end{pmatrix}$$
\end{ex}

\begin{ex}[Moulin.06]{}
	Soit $A\in\mathcal{M}_n(\C)$ et $B$ la transposée de la comatrice de $A$.
	\begin{enumerate}
		\item Montrer que tout vecteur propre de $A$ est vecteur propre de $B$.
		\item On suppose désormais, et pour toute la suite de l'exercice, que $A$ est diagonalisable. Si $A$ est inversible, quelles sont les valeurs propres de $B$ ?
		\item Si $A$ n'est pas inversible, quelles sont les valeurs propres de $B$ ?
	\end{enumerate}
\end{ex}

\begin{ex}{Une forme linéaire pour les valeurs propres ?}
	Existe-t-il une forme linéaire $\Phi$ sur $\mathcal{M}_n(\C)$ telle que pour tout $A\in\mathcal{M}_n(\C)$, $\Phi(A)$ soit une valeur propre de $A$ ?
\end{ex}
\begin{proof}
	Regarder les matrices élémentaires. Cette forme linéaire doit valoir $0$ sur toutes les matrices élémentaires "non diagonales", et vaut $0$ ou $1$ sur les matrices élémentaires "diagonales". Regarder ensuite $I_n$ pour voir qu'exactement une matrice élémentaire "diagonale" a pour image $1$, et les autres $0$. Chercher ensuite une matrice dont le coefficient correspondant n'est pas une valeur propre.
\end{proof}

\begin{ex}{Norme d'algèbre et rayon spectral}
	Soit $N$ une norme d'algèbre sur $\mathcal{M}_n(\C)$ et $\rho(A)$ le rayon spectral de la matrice $A$. Alors, $\rho(A)\leqslant N(A)$.
\end{ex}

\section{Polynôme caractéristique}

\begin{ex}[Moulin.07]{}
	Quels sont les coefficients des termes de degrés $n-2$ et $1$ du polynôme caractéristique d'une matrice de taille $n$ ?
\end{ex}

\begin{ex}[Moulin.08]{$\star$}
	Montrer que tout plan vectoriel de $\mathcal{M}_n(\C)$ ($n\geqslant 2$) contient au moins une matrice singulière non nulle (une matrice singulière est une matrice non inversible).
\end{ex}
\begin{proof}
	On considère $(A,B)$ une base d'un tel plan. Si $B$ est singulière, OK. On peut donc supposer $B$ inversible. Alors : $$\forall\lambda\in\C,\ \det(A=\lambda B)=\underbrace{\det(B)}_{\ne 0}\underbrace{\det(\lambda I_n-B\inv A)}_{\text{ de degré }n}$$ D'après le théorème de d'Alembert-Gauss, il existe donc $\lambda\in\C$ tel que $\det(A+\lambda B)=0$. Et $A+\lambda B\ne 0$ car $1\ne 0$ et $(A,B)$ est libre. Donc $A+\lambda B$ est bien une matrice singulière non nulle.
\end{proof}

\begin{ex}[Moulin.09]{}
	Soit $N\in\mathcal{M}_n(\K)$ une matrice nilpotente. Que dire de $\det(A+N)$ pour $A\in\mathcal{GL}_n(\K)$ commutant avec $N$ ?
\end{ex}
\begin{proof}
	On commence par le cas simple où $A=I_n$. En trigonalisant $N$, $I_n+N$ est semblable à une matrice triangulaire supérieure avec des 1 sur la diagonale donc $$\det(I_n+N)=1=\det(I_n)$$ Dans le cas général, on écrit $$\det(A+N)=\det(A)\det(I_n+A\inv N)$$ Or, $A$ et $N$ commutent, donc $A\inv$ et $N$ commutent et $A\inv N$ est nilpotente. D'après le cas précédent, on en déduit que $$\det(A+N)=\det(A)\times 1=\det(A)$$ 
\end{proof}

\begin{ex}[Moulin.16]{$\det(A+N)$ avec commutation et nilpotence}
	Soit $N\in\mathcal{M}_n(\K)$ une matrice nilpotente. Montrer que si $A\in\mathcal{M}_n(\K)$ commute avec $N$, alors $\det(A)=0\implies\det(A+N)=0$ (voir l'exercice précédent).
\end{ex}
\begin{proof}[Première méthode]
	Raisonnement par densité avec $A_\lambda=A-I\lambda$ qui est inversible au voisinage de $0$ et commute avec $N$. On utilise l'exercice précédent et on conclut par continuité du déterminant ou polynomialité.
\end{proof}
\begin{proof}[Deuxième méthode (sans utiliser l'exercice précédent)]
	On considère $a$ et $n$ les endomorphismes canoniquement associés. Alors par commutation $\ker(a)$ est stable par $n$. Dans une base adaptée à $\ker(a)$ : $$A=\begin{pmatrix}
		(0)&(*)\\
		(0)&(*)
	\end{pmatrix}\ \ \text{et} \ \ N=\begin{pmatrix}
		P&(*)\\
		(0)&(*)
	\end{pmatrix}$$ Donc $$A+N=\begin{pmatrix}
		P&(*)\\
		(0)&(*)
	\end{pmatrix}$$ Or, l'induit de $n$ sur $\ker(a)$ est aussi nilpotent et $\ker(a)$ n'est pas l'espace nul donc $P$ est non inversible. Un déterminant triangulaire par blocs conclut.
\end{proof}

\begin{ex}[Moulin.10]{$\Delta$}
	\begin{enumerate}
		\item $\star$ Soit $(A,B)\in\mathcal{M}_n(\K)^2$. Montrer que $AB$ et $BA$ ont le même polynôme caractéristique. On pourra suppose que $\K$ est un sous-corps de $\C$ et commencer par le cas où l'une des deux matrices est inversible.
		\item Soit $A\in\mathcal{M}_{n,p}(\K)$ et $B\in\mathcal{M}_{p,n}(\K)$. Trouver une relation simple entre les polynômes caractéristiques de $AB$ et $BA$. On pourra remplacer $A$ par une matrice équivalente plus simple. Retrouver alors le résultat de la première question.
		\item Retrouver le résultat de la deuxième question à partir de la première question en complétant les matrices $A$ et $B$ avec des $0$ pour les rendre carrées.
	\end{enumerate}
\end{ex}

\begin{ex}[Moulin.11]{$\Delta$}
	Soit $u$ un endomorphisme d'un espace vectoriel de dimension finie, et $\lambda$ une valeur propre de $u$ d'ordre $p$. Montrer que les conditions suivantes sont équivalentes :
	\begin{enumerate}[label=(\roman*)]
		\item $p=\dim\ker(u-\lambda\text{id}_E)$ ;
		\item $E=\ker(u-\lambda\text{id}_E)\oplus\text{Im}(u-\lambda\text{id}_E)$ ;
		\item $\ker(u-\lambda\text{id}_E)$ possède un supplémentaire stable par $u$.
	\end{enumerate}
\end{ex}

\begin{ex}[Moulin.12]{$\Delta$ Trace et nilpotence}
	\begin{enumerate}
		\item Montrer que $A\in\mathcal{M}_n(\C)$ est nilpotente si, et seulement si, pour tout $k\in\llbracket1,n\rrbracket$, $\Tr(A^k)=0$.
		\item Montrer que $A$ et $B$ dans $\mathcal{M}_n(\C)$ ont même polynôme caractéristique si, et seulement si, pour tout $k\in\N$, $\Tr(A^k)=\Tr(B^k)$.
		\item $\star\star$ Montrer, plus généralement, que $A$ et $B$ dans $\mathcal{M}_n(\C)$ ont même polynôme caractéristique si, et seulement si, pour tout $k\in\llbracket1,n\rrbracket$, $\Tr(A^k)=\Tr(B^k)$.
		\item Application : montrer qu'une matrice $A\in\mathcal{M}_n(\C)$ telle que $\forall k\in\llbracket1,n-1\rrbracket$, $\Tr(A^k)=0$ et $\Tr(A^n)\ne 0$ est nécessairement diagonalisable.
	\end{enumerate}
\end{ex}
\begin{proof}
	\begin{enumerate}
		\item Dans le sens direct, trigonaliser $A$ : toutes ses valeurs propres sont nulles et les traces des puissances valent des sommes de $0$ donc $0$. Dans le sens réciproque, on trigonalise $A$ puis on fait apparaître une matrice de Vandermonde avec les puissances des valeurs propres distinctes, qui multipliée par le vecteur des valeurs propres multipliées par leurs multiplictés fait $0$. On en déduit que la seule valeur propre est $0$ ce qui conclut.
		\item Dans le sens direct, $A$ et $B$ ont donc les même valeurs propres. On trigonalise et le résultat s'obtient facilement. Dans le sens réciproque, on note $\lambda_i$ et $\mu_i$ les valeurs propres respectives de $A$ et $B$ comptées avec multiplicité. Par linéarité, on a l'égalité : $$\forall P\in\C[X],\ \sum_{i=1}^{n}P(\lambda_i)=\sum_{i=1}^{n}P(\mu_i)$$ car cette relation est linéaire en $P$ et on a l'égalité sur la base canonique de $\C[X]$ en traduisant l'hypothèse après trigonalisation. En utilisant les polynômes interpolateurs de Lagrange associés à la réunion des spectre de $A$ et $B$, on trouve que tout élément de cette réunion à même multiplicité chez $A$ et chez $B$. Donc $A$ et $B$ ont même polynôme caractéristique.
		\item Le sens direct est toujours immédiat. Dans le sens réciproque, on cette fois seulement : $$\forall P\in\C_n[X],\ \sum_{i=1}^{n}P(\lambda_i)=\sum_{i=1}^{n}P(\mu_i)$$ et on veut montrer que $$\prod_{i=1}^{n}(X-\lambda_i)=\prod_{i=1}^{n}(X-\mu_i)$$ ce qui revient à une égalité des fonctions symétriques élémentaires. Or, on a l'égalité des sommes de Newton. Il suffit alors de prouver que les sommes de Newton et les fonctions symétriques élémentaires s'expriment les unes en fonction des autres. Cela se fait, mais c'est long et difficile. Se référer au premier DM de l'année.
		\item On pose $x=\Tr(A^n)\ne 0$. On considère alors : $$B=\frac{x}{\sqrt[n]{n}}\text{Diag}(\omega,\ldots,\omega^{n-1})$$ avec $\omega=e^{i2\pi/n}$. Alors les traces des $k$ premières puissances de $B$ sont égales aux traces de $k$ premières puissances de $A$. Donc $A$ et $B$ ont même polynôme caractéristique. Or, $B$ est diagonalisable, donc $A$ aussi. Autre façon de faire : $X^n-\gamma^n=0$. On considère $C$ la matrice compagnon de ce polynôme, et $C$ convient (ses puissances décalent les diagonales).
	\end{enumerate}
\end{proof}

\begin{ex}[Moulin.13]{$\star$}
	Soit $(A,B)\in\mathcal{M}_n(\C)^2$ tel que $A$ et $AB-BA$ commutent. Montrer en utilisant l'exercice précédent que $AB-BA$ est nilpotente.
\end{ex}
\begin{proof}
	Il suffit de montrer que les traces des puissances sont nulles. Or pour $k\in\N$ : \begin{align*}
		(AB-BA)^{k+1}&=(AB-BA)(AB-BA)^k\\
		&=AB(AB-BA)^k-BA(AB-BA)^k \\
		&=A\left[B(AB-BA)^k\right]-\left[B(AB-BA)^k\right]A
	\end{align*} par commutation. Donc $\Tr((AB-BA)^{k+1})=0$. Donc $AB-BA$ est nilpotente.
\end{proof}

\begin{ex}[Moulin.14]{$\Delta$}
	Soit $A$ et $B$ dans $\mathcal{M}_n(\K)$. On suppose qu'il existe $X$ dans $\mathcal{M}_n(\K)$ de rang $r$ tel que $AX=XB$. Montrer que le PGCD de $\chi_A$ et $\chi_B$ est de degré au moins $r$.
\end{ex}
\begin{proof}
	Sr ramener à $X=J_r$ par équivalence. Ensuite, écrire $A$ et $B$ par blocs. La condition donne le fait que les blocs supérieures gauches de $A$ et $B$ sont égaux, que le bloc supérieure droit de $B$ est nul et que le bloc inférieur gauche de $A$ est nul. Donc le polynôme caractéristique du bloc supérieur gauche divise les polynômes caractéristiques de $A$ et de $B$, donc leur PGCD.
\end{proof}

\begin{ex}[Moulin.15]{$\star$}
	Soit $A$ et $B$ deux matrices de $\mathcal{GL}_n(\C)$.
	\begin{enumerate}
		\item On suppose que $\forall M\in\mathcal{M}_n(\C),\ \chi_{AM+B}=\chi{AM}$. Montrer que $B$ est nilpotente. Montrer que $$\forall M\in\mathcal{M}_n(\C),\ \Tr(BAM)=0$$ En conclure que $BA=0$.
		\item Réciproquement, on suppose $B$ nilpotente et $BA=0$. Montrer que $$\forall M\in\mathcal{M}_n(\C),\ \chi_{AM+B}=\chi_{AM}$$
	\end{enumerate}
\end{ex}
\begin{proof}
	\begin{enumerate}
		\item Pour montrer que $B$ est nilpotente, prendre $M=0$ et utiliser le théorème de Cayley-Hamilton. Ensuite, en trigonalisant, on a : $$\forall k\in\N\st, \ \Tr((AM+B)^k)=\Tr((AM)^k)$$ Or, comme $B$ est nilpotente, on a $$\forall k\in\N\st,\ \Tr(B^k)=0$$ Pour $k=2$, on obtient : $$\Tr((AM)^2)+\Tr(AMB)+\Tr(BAM)+\underbrace{\Tr(B^2)}_{=0}=\Tr((AM)^2)$$ donc $2\Tr(BAM)=0$ et on conclut car $\C$ est de caractéristique nulle. Ensuite, c'est l'exercice classique : ici prendre $M=(\overline{BA})^T$ va plus vite.
		\item Réciproquement, on calcule en utilisant le fait pour $\lambda\in\C\st$, $B-\lambda I_n$ est inversible. Alors : \begin{align*}
			\det(\lambda I-B-AM)&=\lambda^n\det(I-(\lambda I-B)\inv AM) \\
			&=\lambda^n\det\left(I-\frac{1}{\lambda}AM\right)\\
			&=\det(\lambda I-AM)
		\end{align*} puisque $(\lambda I-B)=\lambda A$ par hypothèse donc $$(\lambda I-B)\inv A=\frac{1}{\lambda}A$$ On conclut puisque $\C$ est infini.
	\end{enumerate}
\end{proof}

\begin{ex}[Révisions]{$\star$ Matrices de permutation}
	Soit $n\in\N\st$.
	\begin{enumerate}
		\item Définir, pour $\sigma\in\mathcal{S}_n$, une matrice $P_\sigma\in\mathcal{GL}_n(\C)$ telle que l'application $\sigma\mapsto P_\sigma$ soit un morphisme de groupes injectif.
		\item Que représente $\Tr(P_\sigma)$ pour la permutation $\sigma$ ?
		\item Exprimer le polynôme caractéristique de $P_\sigma$ en fonction de la décomposition de $\sigma$ en produit de cycles à supports disjoints.
		\item $\star$ Montrer que deux permutations $\sigma$ et $\tau$ sont conjuguées si, et seulement si, les matrices $P_\sigma$ et $P_\tau$ sont semblables.
	\end{enumerate}
\end{ex}

\section{Sous-espaces stables, commutants}

\begin{ex}[Moulin.17]{}
	Déterminer les sous-espaces vectoriels stables par les endomorphismes canoniquement associés aux matrices suivantes : 
	\begin{enumerate}
		\item $A=\displaystyle\begin{pmatrix}
			2&1&1\\
			1&2&1\\
			0&0&3
		\end{pmatrix}$ ;
		\item $B=\displaystyle\begin{pmatrix}
			1&j&j^2\\
			j&j^2&1\\
			j^2&1&j
		\end{pmatrix}$ ;
		\item $C=\displaystyle\begin{pmatrix}
			1&2&3\\
			0&4&5\\
			0&0&6
		\end{pmatrix}$ ;
		\item $D=\displaystyle\begin{pmatrix}
			2&1&1\\
			1&1&1\\
			-1&-1&0
		\end{pmatrix}$.
	\end{enumerate}
\end{ex}

\begin{ex}[Moulin.18]{}
	En dimension finie égale à $n$, montrer que si $u$ est nilpotent d’indice $n$, alors les sous-espaces stables par $u$ sont exactement les $\ker(u^k)$ pour $k\in\llbracket0, n\rrbracket$.
\end{ex}
\begin{proof}
	En effet, soit $F$ un sous-espace stable par $u$ de dimension $k$. Alors $\chi_{u_F}=X^k$ donc $u_F^k=0$ donc $F\subset\ker(u^k)$, et on a l'égalité par dimension (en utilisant la concavité des noyaux itérés).
\end{proof}

\begin{ex}[Moulin.19]{}
	Quels sont les sous-espaces vectoriels stables par l'endomorphisme : $$\begin{array}{rcl}
		\K[X]&\longrightarrow&\K[X]\\
		P&\longmapsto&P(X+1)-P(X)
	\end{array}\ ?$$
\end{ex}

\begin{ex}[Moulin.20]{$\Delta$}
	Déterminer le commutant d'une matrice diagonalisable.
\end{ex}

\begin{ex}[Moulin.21]{$\star$ Endomorphismes cycliques}
	Soit $u$ un endomorphisme d'un espace vectoriel $E$ de dimension finie $n$. On suppose que $u$ est cyclique, \textit{ie} qu'il existe un vecteur $x$ de $E$ tel que $$E=\text{Vect}(x,u(x),\ldots,u^{n-1}(x))$$ Montrer que les sous-espaces vectoriels de $E$ stables par $u$ sont en nombre fini et tous de la forme $\text{Vect}\{u^k(y)\mid k\in\N\}$ avec $y\in E$.
\end{ex}
\begin{proof}
	S'intéresser à des polynômes annulateurs.
\end{proof}

\begin{ex}[Châteaux]{$\star$ CNS de cyclicité}
	Soit $E$ un $\K$-espace vectoriel de dimension finie et $f$ un endomorphisme de $E$. Démontrer qu'il existe un nombre fini de sous-espaces vectoriels de $E$ stables par $f$ si, et seulement si, $f$ est cyclique.
\end{ex}

\begin{ex}[Moulin.22]{$\Delta$}
	Soit $A_1,\ldots,A_n$ des matrices nilpotentes de $\mathcal{M}_n(\K)$ qui commutent deux à deux. Montre que $A_1A_1\cdots A_n=0$.
\end{ex}
\begin{proof}
	Les noyaux sont stables par commutation. La nilpotence diminue strictement le rang à chaque nouvelle matrice. il est plus facile de le formaliser en termes d'applications linéaires (sinon, cela est rapide en les cotrigonalisant).
\end{proof}

\begin{ex}[Châteaux]{$\star$ Bicommutant d'une matrice diagonalisable}
	Pour $A\in\mathcal{M}_n(\C)$, on note $\mathcal{C}(A)$ le commutant de $A$. On pose $$\mathcal{C}_2(A)=\left\{M\in\mathcal{M}_n(\C)\ \mid\ \forall N\in\mathcal{C}(A),\ MN=NM)\right\}$$ le bicommutant de $A$. Démontrer que si $A$ est diagonalisable, alors : $$\mathcal{C}_2(A)=\C[A]$$
\end{ex}
\begin{rmq}
	Ce résultat est vrai dans le cas général, mais sa démonstration est alors plus difficile et nécessite la décomposition de Frobenius.
\end{rmq}

\section{Diagonalisation}

\begin{ex}[Moulin.23]{$\circ$}
	Diagonaliser les matrices suivantes : 
	$$\begin{pmatrix}
		3&-1&1\\
		-2&4&2\\
		-1&1&5
	\end{pmatrix}\ \ \ \
	\begin{pmatrix}
		1&2&2\\
		2&1&2\\
		2&2&3
	\end{pmatrix}\ \ \ \
	\begin{pmatrix}
		7&2&-2\\
		2&4&-1\\
		-2&-1&4
	\end{pmatrix}$$
\end{ex}

\begin{ex}[Moulin.24]{}
	Diagonaliser $\displaystyle\begin{pmatrix}
		(a)&&b\\
		&\iddots&\\
		b&&(a)
	\end{pmatrix}$ avec $(a,b)\in\R^2$.
\end{ex}
\begin{proof}
	Les lignes sont de somme constante !
\end{proof}

\begin{ex}[Moulin.25]{}
	Calculer les puissances $n$-ièmes $(n\in\N$ puis $n\in\Z$ si cela est possible) des matrices proposées par chacune des méthodes suivantes. Comparer.
	\begin{itemize}
		\item En diagonalisant complètement.
		\item En utilisant un polynôme annulateur.
	\end{itemize}
	$$A=\begin{pmatrix}
		1&-1\\
		2&4
	\end{pmatrix}\ \ \ B=\begin{pmatrix}
	2&0&0\\
	-2&3&1\\
	0&0&2
	\end{pmatrix}\ \ \text{et}\ \ C=\begin{pmatrix}
	1&0&-1\\
	0&1&0\\
	-1&2&1
	\end{pmatrix}$$
\end{ex}

\begin{ex}[Moulin.26]{$\circ$}
	La matrice $A=\displaystyle\begin{pmatrix}
		0&1&2&3\\
		0&4&5&6\\
		0&0&7&8\\
		0&0&0&9
	\end{pmatrix}$ est-elle diagonalisable ? A quelle condition la matrice $B=\displaystyle\begin{pmatrix}
		\lambda I_p&X\\
		(0)&\mu I_q
	\end{pmatrix}$ avec $X\in\mathcal{M}_{p,q}(\K)$ et $(\lambda,\mu)\in\K^2$ est-elle diagonalisable ?
\end{ex}
\begin{proof}
	\begin{itemize}
		\item Elle est triangulaire supérieure donc son polynôme caractéristique est $$\chi_A=X(X-4)(X-7)(X-9)$$ Il est scindé à racines simples donc $A$ est diagonalisable.
		\item Si $\lambda=\mu$, $B$ est diagonalisable si, et seulement si, $X=0$. Si $\lambda\ne \mu$, $B$ est toujours diagonalisable (passer par la dimension des sous-espaces propres).
	\end{itemize}
\end{proof}

\begin{ex}[Moulin.27]{}
	Soit $(a,b,c,d)\in\K^4$ et $M=\displaystyle\begin{pmatrix}
		a&0&0&0\\
		1&b&0&0\\
		0&1&c&0\\
		0&0&1&d
	\end{pmatrix}\in\mathcal{M}_4(\K)$. Montrer que $M$ est diagonalisable si, et seulement si, $a$, $b$, $c$ et $d$ sont deux à deux distincts.
\end{ex}

\begin{ex}[Moulin.28]{$\Delta$}
	On considère $$A=\begin{pmatrix}
		0&1&&(0)\\
		1&0&\ddots&\\
		&\ddots&\ddots&1 \\
		(0)&&1&0
	\end{pmatrix}$$
	\begin{enumerate}
		\item Déterminer les valeurs propres et les vecteurs propres de $A$ en utilisant une suite récurrente. On pourra commencer par chercher les valeurs propres réelles en les mettant sous la forme $2\cos(\theta)$ ou $\pm2\cosh(\theta)$.
		\item Déterminer une matrice de passage $P$ convenable et son inverse (on pourra calculer $P^TP$).
	\end{enumerate}
\end{ex}
\begin{proof} Remarquons que la matrice est symétrique réelle donc orthodiagonalisable : ainsi, toutes ses valeurs propres sont réelles et on peut trouver une matrice de passage orthogonale, ce qui explique les indications.
	\begin{enumerate}
		\item L'équation aux éléments propres fournit la relation de récurrence linéaire suivante pour $AX=\lambda X$ : $$\forall k\in\llbracket1,n\rrbracket,\ x_{k-1}+x_{k+1}=\lambda x_k$$ Si $\lvert\lambda\rvert <2$, on écrit $\lambda=2\cos(\theta)$ et $$\forall x\in\llbracket1,n\rrbracket,\ x_k=A\cos(k\theta)+B\sin(k\theta)$$ Or, $x_0=A=0$ et $$x_{n+1}=\underbrace{B}_{\ne0}\sin((n+1)\theta)=0$$ donc $$\theta=\frac{l\pi}{n+1}$$ On a donc les vecteurs propres $$X_k^{(l)}=\begin{pmatrix}
			\sin\left(\frac{l\pi}{n+1}\right)\\
			\vdots \\
			\sin\left(\frac{ln\pi}{n+1}\right)
		\end{pmatrix}$$ et la matrice de passage $$P=\left(\sin\left(\frac{kl\pi}{n+1}\right)\right)_{1\leqslant k,l\leqslant n}$$ Ensuite, si $\lvert\lambda\rvert >2$, on a une matrice à diagonale domianante donc inversible donc pas de valeur propre. Mais de toute façon, on a déjà $n$ valeurs propres distinctes donc il n'y en a pas d'autre.
		\item En calculant (formules de trigo et sommes géométriques d'exponentielles), on trouve $P^TP=\displaystyle\frac{n}{2}I_n$. On prend finalement $$\tilde{P}=\sqrt{\frac{2}{n}}P\in\mathcal{O}_n(\R)$$ comme matrice de passage.
	\end{enumerate}
\end{proof}

\begin{ex}[Moulin.29]{$\Delta$ Matrices de rang $1$}
	Soit $n\geqslant 2$ et $A\in\mathcal{M}_n(\K)$ de rang $1$.
	\begin{enumerate}
		\item Montrer qu'il existe $X$ et $Y$ des éléments non nuls de $\K^n$ tels que $A=XY^T$.
		\item En déduire un polynôme annulateur de $A$ de degré $2$.
		\item Donner une condition nécessaire et suffisante sur la trace de $A$ pour que $A$ soit diagonalisable.
		\item Déterminer, à similitude près, toutes les matrices de rang $1$ et caractériser les classes de similitude de ces matrices en fonction de leur trace.
	\end{enumerate}
\end{ex}

\begin{ex}[Moulin.30]{$\star$}
	Résoudre dans $\mathcal{M}_3(\R)$ puis dans $\mathcal{M}_3(\C)$ l'équation : $$X^2=\begin{pmatrix}
		-1&-1&-4\\
		-4&5&2\\
		-10&8&-4
	\end{pmatrix}$$
\end{ex}

\begin{ex}[Moulin.31]{}
	Résoudre dans $\mathcal{M}_2(\R)$ puis dans $\mathcal{M}_2(\C)$ l'équation :
	$$M^2+M+I_2=\begin{pmatrix}
		0&1\\
		1&0
	\end{pmatrix}$$
\end{ex}

\begin{ex}[Moulin.32]{}
	Soit $u$ un endomorphisme d'un espace vectoriel de dimension finie. Montrer que $u$ est diagonalisable si, et seulement si, $\chi_u$ est scindé et tout sous-espace propre de $u$ possède un supplémentaire stable par $u$.
\end{ex}

\begin{ex}[Moulin.33]{$\Delta$}
	Soit $u$ un endomorphisme d'un $\C$-espace vectoriel de dimension finie. Montrer que $u$ est diagonalisable si, et seulement si, tout sous-espace vectoriel stable par $u$ admet un supplémentaire stable par $u$.
\end{ex}
\begin{proof}
	Dans le sens direct, si $u$ est diagonalisable et $F$ est stable par $u$, alors $u_F$ est diagonalisable. on complète une base de diagonalisation de$u_F$ avec des vecteurs propres de $u$ pour obtenir un supplémentaire stable par $u$. Réciproquement, supposons que tout SEV stable par $u$ admette un supplémentaire stable par $u$. On pose $$F=\bigoplus_{\lambda\in \text{Sp}(u)}E_\lambda(u)$$ qui est un SEV stable par $u$. Alors il admet $G$ un supplémentaire stable par $u$. Or, il ne peut y avoir de vecteurs propres dans $G$ car ils sont tous dans $F$. Donc puisqu'on travaille sur $\C$, $G=\{0\}$. Donc $$E=F=\bigoplus_{\lambda\in \text{Sp}(u)}E_\lambda(u)$$ et $u$ est bien diagonalisable.
\end{proof}

\begin{ex}[Moulin.34]{}
	Soit $u$ un endomorphisme d'un $\K$-espace vectoriel de dimension finie. Montrer que $u$ est diagonalisable si, et seulement si, tout sous-espace vectoriel admet un supplémentaire stable par $u$. On pourra utiliser des hyperplans.
\end{ex}

\begin{ex}[Moulin.35]{}
	Montrer qu'un endomorphisme d'un espace vectoriel de dimension finie est diagonalisable si, et seulement s'il est combinaison linéaire de projecteurs $(p_i)_{0\leqslant i\leqslant k}$ tels que : $$\forall (i,j)\in\llbracket1,k\rrbracket^2,\ i\ne j\implies p_i\circ p_j=0$$ Quels sont alors son spectre et ses sous-espaces propres ?
\end{ex}

\begin{ex}[Moulin.36]{}
	Déterminer les $A\in\mathcal{M}_n(\C)$ telles que pour tout $P\in\mathcal{GL}_n(\C)$, la matrice $PA$ soit diagonalisable.
\end{ex}

\begin{ex}[Moulin.37]{}
	Soit $A$, $B$ et $M$ dans $\mathcal{M}_n(\C)$. On suppose $A$ et $B$ diagonalisables, et on suppose qu'il existe $r\in\N$ avec $r\geqslant 2$ tel que $A^rMB^r=0$. Montrer que $AMB=0$.
\end{ex}

\begin{ex}[Moulin.38]{$\star$}
	Donner une condition sur $(a_1,\ldots,a_n)\in\C^n$ et $(b_1,\ldots,b_n)\in\C^n$ pour que la matrice $$A=\begin{pmatrix}
		0&a_1&\cdots&a_n\\
		b_1&&&\\
		\vdots&&(0)&\\
		b_n&&&
	\end{pmatrix}$$ soit diagonalisable. On pourra introduire, sauf dans un cas particulier traité à part, un plan vectoriel simple incluant l'image, et examiner l'endomorphisme induit sur ce plan vectoriel.
\end{ex}

\begin{ex}[Moulin.39]{$\Delta$}
	Soit $M\in\mathcal{M}_n(\C)$ et $p\in\N\st$. Montrer que $M$ est diagonalisable si, et seulement si, $M^p$ est diagonalisable et $\ker(M)=\ker(M^p)$.
\end{ex}
\begin{proof}
	Le sens direct est immédiat. Dans le sens réciproque, s'intéresser à des supplémentaires du noyau. Sur celui-ci; utiliser un polynôme annulateur
\end{proof}

\begin{ex}[Moulin.40]{}
	Soit $A\in\mathcal{M}_{4,3}(\R)$ et $B\in\mathcal{M}_{3,4}(\R)$ telles que $BA=\displaystyle\begin{pmatrix}
		0&1&1\\
		1&0&1\\
		1&1&0
	\end{pmatrix}$. La matrice $AB$ est-elle diagonalisable ?
\end{ex}

\begin{ex}{$\star$ Une extension algébrique de $\Q$}
	On pose $E=\Q[j,\sqrt[3]{2}]$ la plus petite $\Q$-algèbre contenant $\{j,\sqrt[3]{2}\}$.
	\begin{enumerate}
		\item Démontrer que $E$ est de dimension $6$ sur $\Q$ et en donner une base.
		\item On considère $\varphi :x\mapsto jx$. Montrer que $\varphi$ est un endomorphisme de $E$.
		\item Déterminer le polynôme caractéristique et le polynôme minimal de $\varphi$.
		\item $\varphi$ est-il diagonalisable ?
	\end{enumerate}
\end{ex}

\begin{ex}[Châteaux]{$\star$}
	On pose $\mathcal{GL}_n(\Z)=\left\{M\in\mathcal{GL}_n(\R)\cap\mathcal{M}_n(\Z)\ :\ M\inv\in\mathcal{M}_n(\Z)\right\}$.
	\begin{enumerate}
		\item Montrer d'une matrice $M\in\mathcal{M}_n(\Z)$ appartient à $\mathcal{GL}_n(\Z)$ si, et seulement si, $\det(M)=\pm1$.
		\item Soit $p$ un nombre premier supérieur ou égal à $3$. On note $\phi$ la projection canonique de $\mathcal{GL}_n(\Z)$ dans $\mathcal{GL}_n(\Z/p\Z)$. Montrer que $\phi$ est un morphisme et que sa restriction a tout sous-groupe fini de l'ensemble de départ est injective.
		\item Peut-on en déduire une majoration sur le cardinal de tels sous-groupes finis ?
	\end{enumerate}
\end{ex}

\begin{ex}[Châteaux]{$\star$}
	Soit $A\in\mathcal{M}_2(\Z)$. On suppose qu'il existe $n\in\N\st$ tel que $A^n=I_2$.
	\begin{enumerate}
		\item Montrer que $A$ est diagonalisable et que ses valeurs propres sont de module $1$.
		\item Quelles sont les valeurs possibles pour le déterminant et la trace de $A$ ?
		\item Soit $G$ un sous-groupe abélien fini de $\mathcal{GL}_2(\Z)$. Montrer que $\forall A\in G,\  A^{12}=I_2$. En déduire que $\card(G)\leqslant 12$.
	\end{enumerate}
\end{ex}

\begin{ex}[Châteaux]{$\star$}
	Soit $A$ et $B$ deux matrices diagonalisables de $\mathcal{M}_2(\C)$. Montrer l'équivalence : $$AB=BA\iff \forall t\in\C,\ A+tB\ \text{est diagonalisable}$$
\end{ex}
\begin{rmq}
	Le résultat reste vrai dans $\mathcal{M}_n(\C)$ et constitue le difficile théorème de Motzkin.
\end{rmq}

\section{Trigonalisation, nilpotents}

\begin{ex}[Moulin.41]{$\circ$}
	Trigonaliser les matrices suivantes :
	$$\begin{pmatrix}
		3&-1&1\\
		-3&5&5\\
		-4&4&6
	\end{pmatrix}\ \ \ \
	\begin{pmatrix}
		3&1&0\\
		-4&-1&0\\
		4&-8&-2
	\end{pmatrix}\ \ \ \
	\begin{pmatrix}
		9&-6&-2\\
		18&-12&-3\\
		18&-9&-6
	\end{pmatrix}$$
\end{ex}

\begin{ex}[Moulin.42]{}
	Soit $M\in\mathcal{M}_n(\C)$ telle que $\text{rg}M=2$, $\Tr M=0$ et $M^n\ne 0$. Montrer que $M$ est diagonalisable.
\end{ex}
\begin{proof}
	Les deux dernières hypothèses impliquent que $M$ possède deux valeurs propres opposées non nulles. Puis avec l'hypothèse sur le rang, on en déduit que tous les sous-espaces propres ont même dimension que la multiplicité de la valeur propre associée, donc que $M$ est diagonalisable.
\end{proof}

\begin{ex}[Moulin.43]{$\star$ Utilisation d'une matrice compagnon}
	\begin{enumerate}
		\item Soit $\displaystyle P=\prod_{k=1}^{n}(X-z_k)\in\Z[X]$ avec $(z_1,\ldots,z_n)\in\C^n$. Montrer que pour tout $N\in\N$, le polynôme $\displaystyle\prod_{k=1}^{n}(X-z_k^N)$ est aussi à coefficients entiers. Indication : considérer la matrice compagnon de $P$.
		\item Application : montrer que si un polynôme à coefficients entiers a toutes ses racines complexes de module inférieur ou égal à 1, alors ses racines non nulles sont des racines de l'unité.
	\end{enumerate}
\end{ex}
\begin{proof}
	\begin{enumerate}
		\item On note $M$ ma matrice compagnon de $P$ qui est à coefficients dans $\Z$. Comme les racines du polynôme sont les valeurs propres de la matrice compagnon, $M$ est semblable à $\tilde{M}$, triangulaire supérieure avec les $z_k$ sur sa diagonale. Par conséquent, $M^N\in\mathcal{M}_n(\Z)$ est semblable à une matrice triangulaire supérieure avec ses coefficients diagonaux égaux aux $z_k^N$. L'égalité des polynômes caractéristiques fournit : $$\prod_{k=1}^{n}(X-z_k^N)=\chi_{M^N}\in\Z[X]$$ donc on a bien le résultat.
		\item Montrons qu'il existe un nombre fini de $P\in\Z[X]$ unitaires de degré $n$ dont toutes les racines sont dans le disque unité. Pour cela, il suffit de borner les fonctions symétriques élémentaires $\sigma_k$ car les coefficients viennent de ces fonctions. Or : $$\lvert\sigma_k\rvert\leqslant\sum_{i_1<\ldots<i_k}\underbrace{\lvert z_{i_1}\cdots z_{i_k}\rvert}_{\leqslant 1}\leqslant\binom{n}{k}$$ Il y a donc un nombre fini de coefficients, donc un nombre fini de polynômes. Soit $P$ correspondant à l'énoncé et $k\in\llbracket1,n\rrbracket$. Alors la suite $(z_k^N)_{N\in\N}$ est à valeurs dans un ensemble fini donc le principe des tiroirs permet de conclure.
	\end{enumerate}
\end{proof}

\begin{ex}[Châteaux]{$\star$ Théorème de Kronecker}
	Soit $A\in\mathcal{M}_n(\Z)$ dont toutes les valeurs propres sont non nulles et de module au plus $1$.
	\begin{enumerate}
		\item Montrer que $\chi_A\in\Z[X]$ et que pour tout $k\in\N$, alors $\Tr(A^k)\in\Z$.
		\item Montrer que les valeurs propres de $A$ sont de module $1$, et que ce sont des racines de l'unité.
		\item Exhiber $A\in\mathcal{M}_n(\Z)$ dont l'ensemble des valeurs propres est $\U_n$.
	\end{enumerate}
\end{ex}

\begin{ex}[Châteaux]{$\star$ $M$ semblable à $2M$}
	Déterminer les $M\in\mathcal{M}_n(\C)$ qui sont semblables à $2M$.
\end{ex}

\begin{ex}[Châteaux]{$\star$}
	Soit $A\in\mathcal{M}_n(\R)$ telle que $A^2$ soit triangulaire supérieure, de diagonale $(1,2,\ldots,n)$. Montrer que $A$ est triangulaire supérieure.
\end{ex}

\begin{ex}[Moulin.44]{$\Delta$}
	Soit $V$ un sous-espace vectoriel de $\mathcal{M}_n(\K)$ de dimension strictement supérieure à $\displaystyle\frac{n(n+1)}{2}$. Montrer que $V$ contient une matrice nilpotente non nulle.
\end{ex}
\begin{proof}
	Considérer l'intersection avec l'ensemble des matrices triangulaires supérieures. Cette intersection n'est pas réduite à l'espace nul car la somme des dimensions est strictement supérieure à $n$. Ainsi, $V$ contient une matrice triangulaire supérieure stricte non nulle, donc une matrice nilpotente non nulle.
\end{proof}
\begin{rmq}
	Dans ce genre d'exercices, toujours considérer une intersection avec un ensemble de matrices de référence possédant la propriété recherchée.
\end{rmq}

\begin{ex}[Moulin.45]{$\star$}
	Soit $E=\mathcal{C}^0([-1,1],\C)$, $F$ un sous-espace vectoriel de $E$ de dimension finie et $f:[-1,1]\to[-1,1]$ croissante et surjective. On suppose que $F$ est stable par l'endomorphisme $D$ de $E$ défini par $$\forall g\in E,\ D(g)=g\circ f$$
	\begin{enumerate}
		\item Montrer que $D$ admet pour seule valeur propre $1$.
		\item En déduire que $D=\text{id}_F$.
	\end{enumerate}
\end{ex}

\begin{ex}[Moulin.46]{}
	Soit $A\in\mathcal{M}_n(\C)$. Quels sont les polynômes $P\in\C[X]$ tels que $P(A)$ soit nilpotente ?
\end{ex}
\begin{proof}
	Trigonaliser. Ce sont ceux dont toutes les valeurs propres de $A$ sont des racines.
\end{proof}

\section{Polynômes d'endomorphismes}

\begin{ex}[Moulin.47]{}
	Condition sur $n\in\N\st$ pour qu'il existe $A\in\mathcal{M}_n(\R)$ telle que $A^2+2A+5I_n=0$ ?
\end{ex}
\begin{proof}
	Un polynôme annulateur a des racines complexes conjuguées non réelles et la matrice est diagonalisable dans $\C$. Les multiplicités en tant que valeurs propres des racines complexes sont égales à celles des conjuguées, donc la dimension est paire. Réciproquement, lorsque la dimension est paire, on construit un bloc $2\times 2$ de la forme : $$A=\begin{pmatrix}
		0&-5\\
		1&-2
	\end{pmatrix}$$ qui convient : on peut le vérifier par le calcul ou plutôt (et c'est mieux) le voir comme la matrice compagnon du polynôme $X^2+2X+5$. On généralise ensuite par blocs.
\end{proof}

\begin{ex}[Moulin.48]{}
	Soit $f$ un endomorphisme de $\R^3$ tel que $f^3+f=0$ et $\text{rg}f=2$. Déterminer les sous-espaces stables par $f$.
\end{ex}

\begin{ex}[Moulin.49]{}
	Soit $A$ une matrice inversible. Montrer que $A\inv$ est un polynôme en $A$.
\end{ex}

\begin{ex}[Moulin.50]{$\star$}
	Soit $u$ un endomorphisme d'un $\K$-espace vectoriel admettant un polynôme minimal $\Pi$. Montrer que si $\Pi$ admet u facteur irréductible $P$ de degré $p$, alors $E$ admet un sous-espace vectoriel de dimension $p$ stable par $u$. On pourra prendre un vecteur non nul de $\ker(P(u))$.
\end{ex}
\begin{proof}
	On écrit $\Pi=P\times Q$ avec $Q(u)\ne0$ (par minimalité de $\Pi$) et $P$ irréductible de degré $p$. On prend alors $x\in E$ non nul tel que $Q(u)(x)\ne0$. On pose alors $y=Q(u)(x)$. On a donc $y\in\ker(P(u))$. On considère alors $$V=\text{Vect}\{u^k(y)\mid k\in\N\}$$ On souhaite montrer que $V$ convient. On a alors $V\subset \ker(P(u))$ car $y\in\ker(P(u))$ et $P(u)$ commute avec tous les $u^k$. Montrons ensuite que la famille $$(y,\ldots,u^{p-1}(y))$$ est une base de $V$. Si on considère une relation de liaison $$\sum_{k=0}^{p-1}\lambda_k u^k(y)=0$$ on pose $R=\displaystyle\sum_{k=0}^{p-1}\lambda_k X^k$. Alors $R$ annule $u_V$. Or, $P$ annule $u_V$ donc $\pi_{u_V}\mid P$. Comme $P$ est irréductible, $\pi_{u_V}$ vaut $1$ ou $P$. Comme $\dim(V)\geqslant1$, $\pi_{u_V}=P$. Mais comme $R$ annule $u_V$ et $\deg(R)<\deg(P)$, on en déduit que $R=0$. Donc la famille $$(y,u(y),\ldots,u^{p-1}(y))$$ est libre. Montrons le caractère générateur. On sait que $$\dim(\K[u_V])=\deg(\pi_{u_V})=\deg(P)=p$$ et on sait aussi que $$V=\K[u_V](y)$$ Il suffit alors de montrer que $$\dim(\K[u_V](y))=p=\dim(\K[u_V])$$ On considère alors l'application $$\begin{array}{lrcl}
		\varphi:&\K[u_V]&\longrightarrow&V\\
		&v&\longmapsto&v(y)
	\end{array}$$ Elle est évidemment surjective. Montrons qu'elle est injective. Si $v(y)=0$, alors pour tout $x\in V$, $v(x)=0$ car $v$ est un polynôme en $u$. Et $v$ coïncide avec l'application nulle sur une base de $V$ donc $v=0$.
\end{proof}

\begin{ex}[Moulin.51]{}
	Soit $A\in\mathcal{M}_3(\R)$ telle que $A^3+A^2+A+I_3=0$ et $A\ne -I_3$. Montrer que $A$ est semblable à $$B=\begin{pmatrix}
		0&-1&0\\
		1&0&0\\
		0&0&1
	\end{pmatrix}$$
\end{ex}
\begin{proof}
	Remarquer que $P=X^3+X^2+X^1+1=(X+1)(X^2+1)$ est un polynôme annulateur de $A$. Donc le spectre de $A$ est inclus dans $\{-1,i,-i\}$. Or, $A$ est de taille 3 et réelle donc admet une valeur propre réelle, qui est nécessairement $-1$. Puis utiliser le lemme des noyaux et remarquer que $\ker(A+I_3)$ ne peut pas être de dimension $3$ comme $A=-I_3$, donc $\ker(A^2+I_3)$ est de dimension $2$. Et on prouve ensuite que la partie correspond à cet espace est semblable à $$\begin{pmatrix}
		0&-1\\
		1&0
	\end{pmatrix}$$
\end{proof}

\begin{ex}[Moulin.52]{$\Delta$}
	Montrer que toute matrice de rang $r$ admet un polynôme annulateur de degré $r+1$.
\end{ex}
\begin{proof}
	On note $u$ l'endomorphisme canoniquement associé à une matrice de rang $r$ et $\chi$ le polynôme annulateur de l'induit de $u$ sur son image. D'après le théorème de Cayley-Hamilton et puisque le rang de la matrice est $r$, $\chi$ est de degré $r$. Puis : $$\forall x,\ \left[\chi \times X\right](u)(x)=\chi(u(x))=0$$ car $u(x)\in\im(u)$. Donc $X\chi$, qui est de degré $r+1$, est un polynôme annulateur de $u$, et donc de la matrice.
\end{proof}

\begin{ex}[Châteaux]{$\star$ Réduction des matrices à coefficients dans $\Z/p\Z$}
	Soit $p$ un nombre premier et $A\in\mathcal{M}_n(\Z/p\Z)$.
	\begin{enumerate}
		\item Montrer que $A$ est diagonalisable si, et seulement si, $A^p=A$.
		\item Démontrer qu'il existe une matrice non trigonalisable dans $\mathcal{M}_2(\Z/p\Z)$.
		\item Démontrer qu'il existe une matrice non trigonalisable dans $\mathcal{M}_n(\Z/p\Z)$.
	\end{enumerate}
\end{ex}

\section{Réduction de matrices blocs}

\begin{ex}[Moulin.53]{}
	Soit $A\in\mathcal{M}_n(\C)$ diagonalisable. Montrer que $$M=\begin{pmatrix}
		A&I_n\\
		I_n&A
	\end{pmatrix}$$ est diagonalisable, et exprimer éléments propres en fonction de ceux de $A$.
\end{ex}

\begin{ex}[Moulin.54]{}
	Soit $A\in\mathcal{M}_n(\R)$. Dans chacun des cas suivants, donner une condition nécessaire et suffisante pour que la matrice blocs proposée soit diagonalisable : $$\begin{pmatrix}
		A&A\\
		1&-3A
	\end{pmatrix}\ \ \text{et}\ \ \begin{pmatrix}
	A&A\\
	-A&-3A
	\end{pmatrix}$$
\end{ex}

\begin{ex}[Moulin.55]{}
	Soit $A\in\mathcal{M}_n(\C)$ et $$B=\begin{pmatrix}
		A&A\\
		0&A
	\end{pmatrix}$$ Donner une condition nécessaire et suffisante sur $A$ pour que $B$ soit diagonalisable.
\end{ex}

\begin{ex}[Moulin.56]{$\star$}
	Condition nécessaire et suffisante sur $A\in\mathcal{M}_n(\R)$ pour que $\displaystyle\begin{pmatrix}
		0&I_n\\
		A&A
	\end{pmatrix}$ soit $\C$-diagonalisable ? $\R$-diagonalisable ?
\end{ex}

\begin{ex}[Moulin.57]{}
	Soit $A\in\mathcal{M}_n(\K)$ et $B\in\mathcal{M}_p(\K)$. 
	\begin{enumerate}
		\item Montrer que si $A$ et $B$ sont diagonalisables, alors $A\otimes B$ l'est également.
		\item Montrer que si $A$ et $B$ sont trigonalisables, alors $A\otimes B$ l'est également.
	\end{enumerate}
\end{ex}

\section{Sous-espaces caractéristiques}

\begin{ex}[Moulin.58]{}
	Soit $u$ un endomorphisme d'un espace vectoriel de dimension finie, et $\lambda$ une valeur propre de $u$ d'ordre $p$. Montrer que les condition suivantes sont équivalentes :
	\begin{enumerate}[label=(\roman*)]
		\item $p=\dim\ker(u-\lambda\text{id}_E)$ ;
		\item $u$ induit sur le sous-espace caractéristique associé à $\lambda$ un endomorphisme diagonalisable ;
		\item $\lambda$ est racine simple du polynôme minimal.
	\end{enumerate}
\end{ex}

\begin{ex}[Moulin.59]{}
	Soit $u$ un endomorphisme d'un $\C$ espace vectoriel de dimension finie $n$. Montrer que les condition suivantes sont équivalentes :
	\begin{enumerate}[label=(\roman*)]
		\item Les sous-espaces propres de $u$ sont tous de dimension $1$.
		\item Il existe un vecteur $x$ de $E$ tel que $(x,u(x),\ldots,u^{n-1}(x))$ soit une base de $E$.
	\end{enumerate}
\end{ex}

\begin{ex}[Moulin.60]{$\star$ Multiplicité dans le polynôme minimal}
	Soit $\lambda$ une valeur propre d'une matrice $A\in\mathcal{M}_n(\R)$. On note $p$ la multiplicité de $\lambda$ en tant que racine du polynôme minimal de $A$, noté $\pi_A$. Montrer que $$\R^n=\ker((A-\lambda I_n)^p)\oplus\im((A-\lambda I_n)^p)$$ 
\end{ex}
\begin{proof}
	On pose : $$P=\frac{\pi_A}{(X-\lambda)^p}$$ Par conséquent, on a $P\wedge(X-\lambda)^p=1$. Comme $\pi_A$ est annulateur de $A$, le lemme des noyaux fournit : $$\R^n=\ker((A-\lambda I_n)^p)\oplus\ker(P(A))$$ Il suffit donc de montrer que $\ker(P(A))=\im((A-\lambda I_n)^p)$. Or : $$\forall X\in\R^n,\ P(A)\times(A-\lambda I_n)^pX=\pi_A(A)X=0\times X=0$$ On a donc l'égalité ensembliste par inclusion (par ce qui précède) et par égalité de dimensions (grâce au théorème du rang et à l'égalité fournie par le lemme des noyaux). On a donc bien :$$\R^n=\ker((A-\lambda I_n)^p)\oplus\im((A-\lambda I_n)^p)$$ 
\end{proof}

\begin{ex}[Moulin.61]{}
	Soit $u$ et $v$ deux endomorphismes commutant d'un $\C$-espace vectoriel $E$ de dimension finie. On suppose que $u$ ne possède aucun sous-espace stable non trivial.
	\begin{enumerate}
		\item En utilisant la décomposition en sous-espace caractéristique, montrer que $v$ est inversible ou nilpotent.
		\item $\star$ Le résultat reste-t-il valable sur un autre corps $\K$ ?
	\end{enumerate}
\end{ex}

\begin{ex}[Moulin.62]{}
	Soit $u$ un endomorphisme d'un $\C$-espace vectoriel de dimension finie. Montrer que $u$ est diagonalisable si, et seulement si : $$\forall P\in\C[X],\ \ker P(u)+\text{Im}P(u)=E$$
\end{ex}

\begin{ex}[Moulin.63]{$\Delta$ Décomposition de Dunford}
	Soit $u$ un endomorphisme d'un espace vectoriel de dimension finie dont le polynôme caractéristique est scindé.
	\begin{enumerate}
		\item Montrer l'existence d'un couple $(d,n)$ d'endomorphismes de $E$, avec $d$ diagonalisable, $n$ nilpotent et $u=d+n$ et $d\circ n=n\circ d$.
		\item Considérons un tel couple $(d,n)$. Soit $\lambda$ une valeur propre de $F_\lambda$ le sous-espace caractéristique associé.
		\begin{itemize}
			\item Vérifier que $u$, $d$ et $n$ stabilisent $F_\lambda$. On note $u_\lambda$, $d_\lambda$ et $n_\lambda$ les endomorphismes qu'ils induisent.
			\item Montrer que $d_\lambda-\lambda\text{Id}_{F_\lambda}$ est nilpotent et en déduire que $d_\lambda=\lambda\text{Id}_{F_\lambda}$.
			\item Montrer l'unicité de $(d,n)$. 
		\end{itemize}
	\end{enumerate}
\end{ex}
\begin{proof}
	A connaitre absolument ! Sert dans un nombre incalculable d'exercices.
	\begin{enumerate}
		\item Utiliser la décomposition en sous-espaces caractéristiques et raisonner matriciellement. On a des blocs diagonaux qui commutent donc par blocs cela commute.
		\item \begin{itemize}
			\item C'est du cours que $u$ le stabilise. $d$ aussi en diagonalisant (à vérifier, je ne suis plus sûr). Et par différence, $n$ aussi.
			\item C'est $-n_\lambda$ qui est aussi nilpotent. Et $d_\lambda-\lambda\text{Id}_{F_\lambda}$ est aussi diagonalisable. Et le seul endomorphisme diagonalisable et nilpotent est l'endomorphisme nul.
			\item On a donc montré que $d$ est déterminé de manière unique par $u$. Donc, par différence, $n$ est aussi déterminé de manière unique par $u$. Donc $(d,n)$ est unique.
		\end{itemize}
	\end{enumerate}
\end{proof}

\begin{ex}[Moulin.64]{$\Delta\star$ Décomposition de Jordan}
	Soit $E$ un $\C$-espace vectoriel de dimension finie non nulle et $u$ un endomorphisme de $E$.
	\begin{enumerate}
		\item On suppose $u$ nilpotent, d'indice de nilpotence $p$.
		\begin{itemize}
			\item Montrer qu'il existe $x\in E$ tel que $\mathcal{B}=(x,u(x),\ldots,u^{p-1}(x))$ soit libre. Montrer que $F=\text{Vect}(\mathcal{B})$ est stable par $u$. Donner la matrice $M_p$, dans la base $\mathcal{B}$, de l'endomorphisme induit par $u$ sur $F$.
			\item Montrer qu'il existe une forme linéaire $\varphi$ sur $E$ telle que $\varphi(u^{p-1}(x))\ne 0$. Montrer que $$G=\bigcap_{k=0}^{p-1}\ker(\varphi\circ u^{k})$$ est stable par $u$ et que $F\oplus G=E$.
			\item Montrer qu'il existe une base de $E$ dans laquelle la matrice de $u$ est diagonale par blocs, avec des blocs de la forme $M_r$.
		\end{itemize}
		\item Montrer qu'il existe une base de $E$ dans laquelle la matrice de $u$ est diagonale par blocs, avec des blocs de la forme $\lambda I_r+M_r$. C'est ce qu'on appelle la forme de Jordan (elle est unique à permutations des blocs diagonaux près).
		\item Application : montrer que toute matrice carrée complexe est semblable à sa transposée.
	\end{enumerate}
\end{ex}
\begin{proof}
	A connaître absolument !
	\begin{enumerate}
		\item \begin{itemize}
			\item Prendre un vecteur tel que $u^{p-1}(x)\ne 0$ (possible car $p$ est l'indice de nilpotence). La famille est alors libre (classique, appliquer $u^{p-1}$ à une relation de liaison, puis $u^{p-1}$, etc.). Vérifier que tous les vecteurs de $\mathcal{B}$ s'envoient sur le suivant, sauf le dernier qui s'envoie sur $0$. On a alors : $$M_p=\begin{pmatrix}
				0&&&(0)\\
				1&0&&\\
				&\ddots&\ddots&\\
				(0)&&1&0
			\end{pmatrix}$$
			\item Prendre la base duale de $\mathcal{B}$ et le dernier vecteur de cette base duale. Idem, vérifier que chaque noyau s'envoie sur le suivant sauf le dernier. Montrer que la somme $F+G$ est directe, ce qui s'obtient facilement par définition de $F$ et $G$. Il est claire que les $\varphi\circ u^k$ sont non nulles (considérer $u^{p-1-k}(x)$), donc $G$ est l'intersection de $p$ hyperplans donc est de dimension supérieure à $n-p$. Donc on a égalité des dimensions et on a bien la supplémentarité (pas besoin d'utiliser l'indépendance des formes linéaires commes ces racistes de MP*2).
			\item Procéder par récurrence en se plaçant sur l'induit de $u$ sur $G$. La récurrence se termine nécessairement.
		\end{itemize}
		\item Utiliser la décomposition en sous-espaces caractéristiques pour avoir diagonalisable plus nilpotent, puis appliquer la première question au nilpotent.
		\item Considérer la forme de Jordan. Ensuite, tout revient à montrer qu'une matrice $M_p$ est semblable à sa transposée. Cela est le cas car cela correspond à un échange des vecteurs de la base.
	\end{enumerate}
\end{proof}

\begin{ex}[LLG]{$\star\star$ Décomposition de Jordan-Dunford multiplicative}
	Soit $E$ un $\K$-ev de dimension finie $n\geqslant1$ et $f\in\mathcal{GL}(E)$. Supposons $\chi_f$ scindé. Montrer qu'il existe un unique couple $(d,u)\in\mathcal{GL}(E)^2$ tel que : \begin{itemize}
		\item $d$ est diagonalisable ;
		\item $u$ est unipotent, \textit{ie} il existe $m\in\N\st$ tel que $(u-\text{id})^m=0$ ;
		\item $f=d\circ u=u\circ d$.
	\end{itemize}
	Montrer de plus que $d$ et $u$ sont des polynômes en $f$.
\end{ex}
\begin{proof}[Existence]
	On écrit la décomposition de $E$ en sous-espaces caractéristiques de $f$ : $$E=\bigoplus_{i=1}^{s}F_{\lambda_i}$$ On note $\pi_i$ la projection sur $F_{\lambda_i}$ parallèlement à la somme des autres sous-espaces caractéristiques. On pose alors $$d=\sum_{i=1}^{s}\lambda_i\pi_i$$ $d$ est inversible, il suffit de considérer une base de trigonalisation associées à la décomposition en sous-espaces caractéristiques, et on voit alors que $d$ est aussi diagonalisable. On pose ensuite : $$u=d\inv\circ f$$ $u$ est unipotente car dans cette base, sa matrice est triangulaire supérieure avec uniquement des $1$ sur la diagonale. Puisque les $\pi_i$ sont des polynômes en $f$, $d$ l'est aussi. Comme $\K[d]$ est une algèbre et $d$ est inversible, $d\inv\in\K[d]$ donc \textit{a fortiori} $d\inv\in\K[f]$, et ainsi $u$ est un polynôme en $f$. Par conséquent, $u$ et $d$ commutent et par construction $f=d\circ u$.
\end{proof}
\begin{proof}[Unicité]
	On considère un couple $(d_1,u_1)$ qui vérifie les mêmes hypothèses. Alors $u$, $u_1$, $d$ et $d_1$ commutent tous entre eux car ce sont des polynômes en $f$. On pose alors $g=u_1\inv u=d_1d\inv$. $g$ est alors diagonalisable car $d_1$ et $d\inv$ sont codiagonalisables car commutent. Puis, $u$ et $u_1\inv$ sont cotrigonalisables donc $g$ est trigonalisable et ses valeurs propres sont les produits des valeurs propres de $u$ et $u_1\inv$, donc $1$. Donc $g$ est unipotente. Comme le seul endomorphisme unipotent et diagonalisable est $\text{id}$ (car $g-\text{id}$ est nilpotent est diagonalisable donc nul), on en déduit que $g=\text{id}$ donc $u=u_1$ et $d=d_1$.
\end{proof}

\section{Coréduction}

\begin{ex}[Moulin.65]{$\Delta$}
	Vrai ou faux ?
	\begin{itemize}
		\item Deux endomorphismes diagonalisables sont codiagonalisables si, et seulement s'ils commutent.
		\item Deux endomorphismes trigonalisables sont cotrigonalisables si, et seulement s'ils commutent.
	\end{itemize}
\end{ex}
\begin{proof}
	\begin{itemize}
		\item Vrai. Le sens réciproque est démontré dans la partie astuces. Pour le sens direct, codiagonaliser et vérifier la commutation sur les sous-espaces de la décomposition.
		\item Faux. Le sens réciproque est vrai et est démontré dans la partie astuces. Mais le sens direct est faux, car deux matrices triangulaires ne commutent pas nécessairement. On pourra prendre par exemple $$\begin{pmatrix}
			2&1\\
			0&1
		\end{pmatrix} \ \ \text{et} \ \ \begin{pmatrix}
			2&2\\
			0&1
		\end{pmatrix}$$
	\end{itemize}
\end{proof}

\begin{ex}[Moulin.66]{$\star$}
	Soit $(A,B)\in\mathcal{M}_n(\K)^2$ un couple de matrices diagonalisables. Montrer que $\varphi :M\mapsto AM-MB$ est diagonalisable. Même question en remplaçant "diagonalisable(s)" par "trigonalisable(s)".
\end{ex}
\begin{proof}
	On pose $\varphi_A:M\mapsto AM$ et $\varphi_B:M\mapsto MB$. On vérifie par récurrence (en mettant à part le cas $k=0$) que $$\varphi_A^k=\varphi_(A^k)$$ Par linéarité, on en déduit que $$\forall P\in\C[X],\ P(\varphi_A)=\varphi_{P(A)}$$ Or, $A$ est annulée par un polynôme scindé à racines simples (car diagonalisable) et alors, comme $\varphi_0=0$, $\varphi_A$ est aussi annulée par un polynôme scindé à racines simples, donc est diagonalisable. De même pour $B$. Or, on voit que $\varphi_A$ et $\varphi_B$ par associativité du produit matriciel. Par conséquent, $\varphi_A$ et $\varphi_B$ sont codiagonalisables, donc $\varphi=\varphi_A-\varphi_B$ est diagonalisable. Pour trigonalisable, il suffit de remplacer dans ce qui précède "diagonalisable" par "trigonalisable" et "scindé à racines simples" par scindé.
\end{proof}

\begin{ex}[Moulin.67]{}
	Soit $f$ et $g$ deux endomorphismes d'un espace vectoriel $E$ de dimension finie. Pour $g\in\mathcal{L}(E)$, on pose $F(g)=g\circ f$ et $H(g)=h\circ g$, définissant ainsi deux endomorphismes de $\mathcal{L}(E)$.
	\begin{enumerate}
		\item Calculer $F^k$ et $H^k$ pour $k\in\N$. En déduire que $F$ est diagonalisable si, et seulement si, $f$ est diagonalisable, et de même pour $H$ et $h$.
		\item Donner une condition suffisante pour que $g\mapsto h\circ g\circ f$ soit diagonalisable.
		\item $\star\star$ On suppose $\K=\C$. Donner une condition nécessaire et suffisante pour que $g\mapsto h\circ g\circ f$ soit diagonalisable.
	\end{enumerate}
\end{ex}

\begin{ex}[Moulin.68]{}
	Reprendre l'exercice précédent en remplaçant "diagonalisable" par "trigonalisable".
\end{ex}

\begin{ex}[Moulin.69]{}
	Soit $$A=\begin{pmatrix}
		1&j&j^1\\
		j&j^2&1\\
		j^2&1&j
	\end{pmatrix}$$ Déterminer les matrices $B\in\mathcal{M}_3(\C)$ telles que $B^2=A$.
\end{ex}

\begin{ex}[Moulin.70]{Une CS de cotrigonalisabilité}
	Soit $A$ et $B$ deux matrices de $\mathcal{M}_n(\C)$ telles que $AB=0$. Montrer que $A$ et $B$ sont simultanément trigonalisables.
\end{ex}
\begin{proof}
	On suppose $B$ non nulle sinon cela n'a pas beaucoup d'intérêt. Alors $\im(B)$ est stable par $B$ donc il existe un vecteur propre de $B$ dans $\im(B)$ car $\im(B)$ n'est pas l'espace nul. L'hypothèse $AB=0$ fournit alors qu'un tel vecteur propre est aussi vecteur propre de $A$ associé à la valeur propre $0$. On a donc un vecteur propre commun à $A$ et $B$. Puis récurrence comme d'habitude.
\end{proof}

\begin{ex}[Châteaux]{Un exemple de cotrigonalisation}
	Soit $E$ de dimension finie $n$ sur $\C$ ainsi que $u$, $v$ et $w$ trois endomorphismes de $E$ vérifiant : $$\begin{cases}
		uv-vu=w\\
		uw=wu\\
		vw=wv
	\end{cases}$$
	\begin{enumerate}
		\item Montrer que $w$ est nilpotent.
		\item Montrer que $u$, $v$ et $w$ ont un vecteur propre commun.
		\item En déduire qu'ils sont trigonalisables dans une même base.
	\end{enumerate}
\end{ex}
\begin{ex}[Moulin.71]{$\Delta$ Crochet de Lie et cotrigonalisation}
	Soit $(A,B)\in\mathcal{M}_n(\C)^2$ tel que $AB-BA=\lambda A+\mu B$ avec $(\lambda,\mu)\in\C^2$. On cherche à démontrer que $A$ et $B$ sont cotrigonalisables.
	\begin{enumerate}
		\item Traiter le cas $\lambda=\mu=0$.
		\item Lorsque $(\lambda,\mu)\ne(0,0)$, montrer que l'on peut se ramener au cas où $(\lambda,\mu)=(1,0)$, ce que l'on supposera dans la suite.
		\item Montrer que $A$ est nilpotente et que $B$ stabilise un sous-espace vectoriel non trivial stable par $A$.
		\item Conclure.
	\end{enumerate}
\end{ex}
\begin{proof}
	\begin{enumerate}
		\item Si $\lambda=\mu=0$, $A$ et $B$ commutent donc sont cotrigonalisables (\textit{cf} démonstration dans la partie astuces).
		\item Sans perte de généralité, quitte à multiplier $AB-BA$ par -1 (ce qui revient à échanger les rôles de $A$ et $B$ mais ne change pas la non nullité du couple de scalaires), on peut supposer $\lambda\ne0$. Puis, sans perte de généralité, quitte à prendre $\tilde{A}=\lambda A$, $\tilde{B}=\frac{1}{\lambda}B$, $\tilde{\mu}=\lambda\mu$ et $\tilde{\lambda}=1$, on peut supposer $\lambda=1$. Enfin, sans perte de généralité, quitte à prendre $\tilde{A}=A+\mu B$ et $\tilde{\mu}=0$, on peut supposer $\mu=0$. En effet, toutes les opérations précédentes n'influent pas sur la cotrigonalisabilité.
		\item On suppose donc désormais $(\lambda,\mu)=(1,0)$. On note : $$\begin{array}{lrcl}
			f_B:&\mathcal{M}_n(\C)&\longrightarrow&\mathcal{M}_n(\C)\\
			&M&\longmapsto&MB-BM
		\end{array}$$ On constate alors que pour $k\in\N$ : \begin{align*}
			f_B(A^{k+1})&=A^{k+1}B-BA^{k+1} \\
			&=A(A^kB-BA^k)+(AB-BA)A^k \\
			&=Af_B(A^k)+A^{k+1}
		\end{align*} On en déduit par une récurrence immédiate : $$\forall k\in\N,\ f_B(A^k)=kA^k$$ Or, on est en dimension finie, donc $f_B$ ne peut avoir qu'un nombre fini de valeurs propres. Par conséquent, $A^k$ est nulle au-delà d'un certain $k$. Donc $A$ est nilpotente. On peut supposer $A$ non nulle et $n\geqslant 2$, sinon le résultat est immédiat. Alors, le noyau de $A$ est stable par $B$ car $AB=BA+A$. Comme $A$ est non nulle et nilpotente, le sous-espace $\ker(A)$ convient.
		\item On a donc une valeur propre commune sur ce sous-espace. puis récurrence comme d'habitude, en le faisant matriciellement et en termes d'endomorphismes en même temps pour éviter de se trimbaler des projections sur des hyperplans. 
	\end{enumerate}
\end{proof}
\begin{rmq}
	On appelle crochet de Lie la quantité : $$[A,B]=AB-BA$$ On peut munir des espaces d'une telle opération pour obtenir la structure d'algèbre de Lie.
\end{rmq}

\begin{ex}[Moulin.72]{}
	Soit $\mathcal{N}$ un sous-espace vectoriel de $\mathcal{M}_n(\K)$ stable par multiplication et constitué de matrices nilpotentes.
	\begin{enumerate}
		\item Soit $X\in\K^n$. Montrer que $\mathcal{N}X=\{NX\ \mid\ N\in\mathcal{N}\}$ est un sous-espace vectoriel de $\K^n$ stable par tout élément de $\mathcal{N}$.
		\item Peut-on avoir $\mathcal{N}X=\K^n$ ?
		\item En déduire qu'il existe une base de trigonalisation simultanée de tous les éléments de $\mathcal{N}$.
	\end{enumerate}
\end{ex}

\begin{ex}[Moulin.73]{}
	Soit $E$ un espace vectoriel de dimension finie et $(f,u)\in\mathcal{L}(E)^2$ avec $u$ inversible. On suppose que $u\circ f+f\circ u=0$.
	\begin{enumerate}
		\item On suppose $f$ diagonalisable. Montrer que $\text{rg}(f)$ est pair.
		\item On suppose $\ker(f)\oplus\text{Im}(f)=E$. Montrer que $\text{rg}(f)$ est pair.
		\item A-t-on $\text{rg}(f)$ pair en général ?
	\end{enumerate}
\end{ex}

\begin{ex}[Moulin.74]{$\star$}
	Soit $(A,B)\in\mathcal{M}_n(\R)$ tel que $AB=BA$, $A^n=B^n=I_n$ et $\Tr(AB)=n$. Montrer que $\Tr(A)=\Tr(B)$.
\end{ex}
\begin{proof}
	On les cotrigonalise dans $\C$. On note $(a_1,\ldots,a_n)$ et $(b_1,\ldots,b_n)$ leurs valeurs propres comptées avec multiplicité. L'hypothèse $A^n=B^n=I_n$ montre que les $a_i$ et les $b_i$ sont des racines n-èmes de l'unité. Par conséquent, les $a_ib_i$ sont des complexes de module $1$. On en déduit que $$\lvert\Tr(AB)\rvert\leqslant\sum_{i=1}^{n}\lvert a_ib_i\rvert=n$$ Le cas d'égalité de l'inégalité triangulaire fournit que tous les $a_ib_i$ sont colinéaires. Et comme tous les $a_i b_i$ sont de module $1$, ils sont tous égaux car si deux sont opposés, la somme ne peut pas faire $n$. Et comme ils sont tous égaux, ils valent tous $1$. Donc $b_i=\overline{a_i}$. On en déduit que $$\Tr(B)=\overline{\Tr(A)}$$ Or, $\Tr(A)\in\R$ puisque $A$ est réelle. Par conséquent, $$\Tr(B)=\Tr(A)$$
\end{proof}

\begin{ex}[ENS]{$\star$}
	\begin{enumerate}
		\item Déterminer une famille libre d'éléments de $\mathcal{M}_n(\C)$ commutant deux à deux et de cardinal $1+\displaystyle\left\lfloor\frac{n^2}{4}\right\rfloor$.
		\item Montrer qu'une famille commutative d'éléments de $\mathcal{M}_n(\C)$ est cotrigonalisable.
		\item Montrer que le cardinal d'une famille libre d'éléments de $\mathcal{M}_n(\C)$ commutant deux à deux est majoré par $1+\displaystyle\left\lfloor\frac{n^2}{4}\right\rfloor$.
	\end{enumerate}
\end{ex}

\begin{ex}[ENS]{$\star$}
	Soit $V$ un $\C$-espace vectoriel de dimension finie. Soit $A$ et $H$ des sous-espaces vectoriels de $\mathcal{L}(V)$ vérifiant : $$\begin{cases}
		\forall (M,N)\in A^2,\ MN-NM\in A\\
		\forall X\in H,\ \forall M\in A,\ XM-MX\in H
	\end{cases}$$
	\begin{enumerate}
		\item Soit $v$ un vecteur propre commun à tous les éléments de $H$. Montrer : $$\forall X\in H,\ \forall M\in A,\ v\in\ker(XM-MX)$$
		Indication : considérer $W=\text{Vect}\left\{M^kv\ \mid\ k\in\N\right\}$ et regarder la matrice d'un élément de $H$ dans une bonne base de $W$.
		\item On suppose que $H$ est un hyperplan de $A$. Soit $F$ un espace propre commun à tous les éléments de $H$. Montrer que $F$ est un espace propre commun à tous les éléments de $A$.
	\end{enumerate}
\end{ex}

\end{document}